% (User input window)
% 2026-02-15 19:56:58（Asia/Singapore）
% According to an internal note dated 2026-02-14.

\begin{corollary}[主指数的双极限刚性：\texorpdfstring{$\tfrac12\log\varphi$}{(1/2)logphi} 为不可混叠的唯一体积斜率]\label{cor:pom-double-limit-rigidity-half-logphi}
对整数 \(q\ge 2\) 与分辨率 \(m\) 令
\(
a_{m,q}:=\frac{1}{mq}\log S_q(m)
\)。
则在已核验的 Perron 主尺度口径下（\(r_q=\lim_{m\to\infty}S_q(m)^{1/m}\) 存在），有
\[
\lim_{q\to\infty}\ \lim_{m\to\infty}\ a_{m,q}
=\lim_{m\to\infty}\ \lim_{q\to\infty}\ a_{m,q}
=\frac12\log\varphi.
\]
\end{corollary}
\begin{proof}
对固定 \(q\)，由 \(S_q(m)=\Theta(r_q^m)\) 得
\(
\lim_{m\to\infty}a_{m,q}=\frac{1}{q}\log r_q
\)。
再由端点定理（命题 \ref{prop:pom-rq-qinfty-endpoint}）有
\(
\lim_{q\to\infty}\frac{1}{q}\log r_q=\frac12\log\varphi
\)，得到第一条迭代极限。
\par
对固定 \(m\)，由最大项主导（\(q\to\infty\) 的 \(\ell^q\) 极限）有
\(
\lim_{q\to\infty}S_q(m)^{1/q}=D_m:=\max_{x\in X_m}d_m(x)
\)，从而
\(
\lim_{q\to\infty}a_{m,q}=\frac{1}{m}\log D_m
\)。
再由最大纤维闭式定理（定理 \ref{thm:pom-max-fiber}）得
\(
\lim_{m\to\infty}\frac{1}{m}\log D_m=\frac12\log\varphi
\)，得到第二条迭代极限。证毕。
\end{proof}

\begin{corollary}[自由能超额控制维数收敛：端点维数的可审计误差分解]\label{cor:pom-Dq-convergence-controlled-by-gq}
令自由能超额 \(g_q\) 如注 \ref{rem:pom-free-energy-excess-gq} 定义：
\[
g_q:=\log r_q-\frac{q}{2}\log\varphi.
\]
则 Rényi 熵率满足恒等式
\[
h_q=q\log\frac{2}{\sqrt{\varphi}}-g_q,
\]
从而在几何标度 \(\varepsilon_m=\varphi^{-m}\) 下，Rényi 维数谱（命题 \ref{prop:pom-renyi-dimension-spectrum}）可写为
\[
\boxed{\;
D_q
=\frac{q}{q-1}\cdot\frac{\log(2/\sqrt{\varphi})}{\log\varphi}
-\frac{g_q}{(q-1)\log\varphi}\;}
\qquad(q\ge 2).
\]
记端点维数
\(
D_\infty:=\frac{\log(2/\sqrt{\varphi})}{\log\varphi}
\)，
则偏差具有严格分解
\[
\boxed{\;
D_q-D_\infty=\frac{1}{q-1}\left(D_\infty-\frac{g_q}{\log\varphi}\right).}
\]
因此任何关于“维数谱是否快速贴近端点”的可证伪陈述，都可被降解为对 \(g_q\) 衰减型态的审计（而 \(g_q\) 的数值证书已由注 \ref{rem:pom-free-energy-excess-gq} 给出）。
\end{corollary}
\begin{proof}
由 \(h_q=q\log 2-\log r_q\) 与 \(g_q=\log r_q-\frac{q}{2}\log\varphi\) 直接整理得
\(
h_q=q\log(2/\sqrt{\varphi})-g_q
\)。
代入 \(D_q=\frac{h_q}{(q-1)\log\varphi}\) 即得第一式；再把常数项记为 \(D_\infty\) 并移项得到偏差分解。证毕。
\end{proof}

